\documentclass[10pt,a4paper]{amsart}
\usepackage[margin=19mm]{geometry}
\usepackage{amsmath,amssymb,mathtools}
\usepackage[scr=rsfs,cal=euler]{mathalfa}
\usepackage{xfrac}
\usepackage[unicode,hidelinks,bookmarks=false]{hyperref}
\newcommand{\ie}{i.e.\ }
\newcommand{\cf}{cf.\ }
\newcommand{\resp}{respectively\ }
\newcommand{\Subsection}{\subsection}
\providecommand{\nobreakdash}{\nobreak\hbox{-}}
\setlength{\emergencystretch}{2em}
\multlinegap0pt
\usepackage{bm}
\usepackage{bbm}
\usepackage{dsfont}
\usepackage{xspace}
\usepackage{cancel}
\usepackage{mathtools}
\usepackage{amsthm}
\makeatletter
\@ifundefined{Lemma}{\newtheorem{Lemma}{Lemma}}{}
\@ifundefined{Theorem}{\newtheorem{Theorem}[Lemma]{Theorem}}{}
\@ifundefined{prop}{\newtheorem{prop}[Lemma]{Proposition}}{}
\makeatother
\newcommand{\kk}{\Bbbk}
\title[Cohen-Macaulay modules of covariants for cyclic $p$-groups]{Cohen-Macaulay modules of covariants for cyclic $p$-groups}
\author{Jonathan Elmer}
\date{Source: arXiv:2506.03677v1 (CC BY 4.0)}
\begin{document}
\maketitle

\noindent\emph{Source and licence.} Cohen-Macaulay modules of covariants for cyclic $p$-groups. Version arXiv:2506.03677v1 (CC BY 4.0), distributed under the Creative Commons Attribution 4.0 International licence, as recorded in the arXiv licence metadata for this exact version and shown on the abstract page https://arxiv.org/abs/2506.03677v1. Full rights notice: licenses/arxiv-2506.03677-CC-BY-4.0.txt.\par\smallskip

\noindent\textit{Editorial scope.} Counted unit: the Theorem whose source label is \texttt{None} in the article (source0.tex lines 586--592, source label \texttt{None}), delivered together with the article's own complete original proof of it (source0.tex lines 595--650, one \texttt{proof} environment of 56 lines). No mathematics is written, repaired or re-derived: statement and proof are the article's own text line for line. Required context retained uncounted: freetest (lines 264--266); v2v2leadterms (lines 569--572). Every definition, display and label of those lines is retained unchanged. Omitted from this package and not redistributed: 1--263, 267--568, 573--585, 593--594, 651--726. Nothing inside a retained excerpt is omitted or altered: every retained body is delivered exactly as the article's lines are written, so no omission receipt of any kind accompanies this package. Citations: the retained excerpts carry no citation command at all, so no bibliography entry of the article is transcribed and no reference is redistributed. Reference adaptation: none was needed and none was made; every reference inside the delivered unit resolves inside it, so no unresolved reference or citation is produced. Printed slips kept literal: no printed word, symbol, hypothesis or number of the counted statement or of its proof was altered. Layout and class: the article's own class is \texttt{amsart} with packages including latexsym, amssymb, amsmath, amsthm, amsfonts; the wrapper is the lane's pinned \texttt{amsart} editorial wrapper at 10pt on A4, which restates the theorem-like environments the retained text needs and adds the article's own macro definitions that the retained text needs (\texttt{\textbackslash kk}), with extra wrapper packages (bm, bbm, dsfont, xspace, cancel, mathtools). The delivered numbering is the unit's own. Assets: no retained span contains a figure environment, a graphics-inclusion command, a PDF-page inclusion, a table or a TikZ picture; no image file is copied into this package, the package declares no asset dependency and no image file is redistributed with it. Licence: version arXiv:2506.03677v1 of this article is distributed under the Creative Commons Attribution 4.0 International licence, as recorded in the arXiv licence metadata for this exact version and shown on the abstract page https://arxiv.org/abs/2506.03677v1. A full rights notice is delivered at licenses/arxiv-2506.03677-CC-BY-4.0.txt. Characters: every retained excerpt is pure ASCII, so the delivered engine, XeLaTeX with its default Latin Modern text font, reports no missing character.
\begin{prop}\label{freetest} Let $A$ be a regular, graded $\kk$-module with $A_0= \kk$ and let $M$ be a finite generated free graded module over $A$. Let $g_1, g_2, \ldots, g_r$ be a $A$-independent set of elements of $M$, where $r = r(M,A)$. Then
\[\sum_{i=1}^r \deg(g_i) \geq s(M,A)\] with equality if and only if $M$ is generated by $g_1, \ldots, g_r$.
\end{prop}

\begin{Lemma}\label{v2v2leadterms}
The lead term of $\Delta^k(x_1^ix_2^j)$ is:
\[\left\{ \begin{array}{lr} \frac{j!}{(j-k)!}x_1^i x_2^{j-k} x_4^k & k \leq j \\ \frac{i! j!}{(i-j-k)!}x_1^{i-k}x_3^{k-j}x_4^j & j<k \leq i+j. \end{array} \right.\]
\end{Lemma}

\begin{Theorem} Let $n \leq p$. The Kernel $K_n$ of $\Delta^n: \kk[V] \rightarrow \kk[V]$ is a free $A$-module generated by the set $S:= M \cup U \cup D$ where
\begin{align*}
M &= \bigcup_{k=0}^{n-1} M_k;\\
U &= \bigcup_{k=1}^{p-n} u^kM_{n-1};\\
D &= \bigcup_{k=2p-n}^{2p-2} \Delta^{p-n}M_k. 
\end{align*}
\end{Theorem}

\begin{proof}
The $A$-isomorphism $K_n \cong \kk[V,W]^G$ shows that $K_n$ is free, and $r(K_n,A) = np$, and $s(K_n,A) = np(p-1)$. 
The elements of $M_k$ for $k \leq n-1$ have weight $k+1$, so the weight of each element of $M$ is at most $n$. Since $u \in \kk[V]^G$, the weight of each element of $U$ is $n$. If $k \geq 2p-n$ then $k \geq p$, so in that case the weight of each element of $M_k$ is $p$. It follows that the weight of each element of $D$ is $n$. The shows that the proposed generators all lie in $K_n$. Now by Proposition \ref{freetest} it is enough to show that the given set is $A$-independent, has $np$ elements, and degree sum $np(p-1)$.

Viewed as a polynomial in $x_1,x_2$ with coefficients in $\kk[x_3,x_4]$, the monomials in $M_k$ each have total degree $k$. Therefore the elements of $M$ have total degree at most $n-1$. Further, each monomial in $M_k$ has a distinct $x_1,x_2$ bidegree. So the elements of $M$ have distinct bidegrees.

Meanwhile, the elements of $u^kM_{n-1}$ have total degree (in $x_1,x_2$) $k+n-1$. Therefore the total degrees of the elements of $U$ lie between $n$ and $p-1$ inclusive. Further, each element of $u^kM_{n-1}$ has distinct bidegree. So the bidegrees of the elements of $U$ are distinct from each other and also from all elements of $M$.

Let $m:= x_1^ix_2^j \in M_k$ with $k \geq 2p-n$. Then $p-n \leq k-p = i+j-p < j$, so by Lemma \ref{v2v2leadterms}, the lead term of $\Delta^{p-n}(m)$ is $$\frac{j!}{(j-p+n)!}x_1^ix_2^{j-p+n}x_4^{p-n}.$$
In particular, its total degree in $x_1,x_2$ is $i+j-p+n = k-p+n \geq p$. So the elements of $D$ have bidegree $\geq p$. Further, the bidegrees of the elements of $\Delta^{p-n}M_k$ are distinct from one another. We have shown that the bidegrees of the elements of $S$ are all distinct from one another.  

Multiplying an element of $\kk[V]$ by an element of $A$ preserves its $x_1$- and $x_2$- degree modulo $p$. Since each element of $S$ has $x_1$- and $x_2$- degree $<p$, and distinct $x_1,x_2$-bidegree, it follows that there is no $A$-linear relation between the elements of $S$. So $S$ is $A$-independent as required.

It remains to check the size and degree sum of $S$. The number of elements in $M_k$ is $k+1$ if $k < p$, and $2p-k-1$ otherwise. Thus
\begin{align*}
|S| &= |M|+|U|+|D|\\
&= \sum_{k=0}^{n-1} (k+1) + n(p-n) + \sum_{k=2p-n}^{2p-2} (2p-1-k)\\
&= \frac12 n(n+1) + n(p-n) + (2p-1)(n-1) - \frac12(n-1)(4p-n-2)\\
&= \frac12 n(n+1) + n(p-n) + \frac12 n(n-1)\\
&= \frac12 n \cdot (2n) + np-n^2\\
&=np
\end{align*}
as required. 

Finally, the degree of each element of $M_k$ is $k$, the degree of each element of $u^kM_{n-1}$ is $2k+n-1$ and the degree of each element of $\Delta^{p-n}M_k$ is also $k$. We note the identity
\[\sum_{k=0}^n k(k+1) = \frac13 n(n+1)(n+2).\]
Therefore the degree sum of $M$ is
\begin{align*}
\sum_{k=0}^{n-1} k(k+1) &= \frac13(n-1)n(n+1)\\
&= \frac13 (n^3-n).
\end{align*}
The degree sum of $U$ is
\begin{align*}
\sum_{k=1}^{p-n} (2k+n-1)n &= 2n \sum_{k=1}^{p-n}k + (p-n)(n-1)n\\
&= n(p-n)(p-n+1)+n(p-n)(n-1)\\
&=np^2-n^2p.
\end{align*}
And the degree sum of $D$ is
\begin{align*}
&\sum_{k=2p-n}^{2p-2} k(2p-1-k)\\
&= 2p \sum_{k=2p-n}^{2p-2}k - \sum_{k=2p-n}^{2p-2}k(k+1)\\
&= p(4p-2-n)(n-1) - \frac13(2p-2)(2p-1)(2p) + \frac13 (2p-n-1)(2p-n)(2p-n+1)\\
&= p(4p-2-n)(n-1) - \frac13(2p-2)(2p-1)(2p)\\
& \qquad + \frac13 (2p-2-(n-1))(2p-1-(n-1))(2p-(n-1))\\
&= p(4p-2-n)(n-1) - \frac13\left((6p-3)(n-1)^2 - (12p^2-12p+2)(n-1)-(n-1)^3 \right)\\
&= (n-1)\left(p(4p-2-n+\frac13(6p-3(n-1)-12p^2+12p-2-(n-1)^2)) \right)\\
&= (n-1)(pn-\frac13n - \frac13n^2 )\\
&= n^2p-np-\frac13(n^3-n).
\end{align*}
So the degree sum of $S$ is
\begin{align*}
&\frac13 (n^3-n)+np^2-n^2p=n^2p-np-\frac13(n^3-n)\\
& = np(p-1)
\end{align*}
as required. This completes the proof.
\end{proof}

\end{document}
